\documentclass[runningheads]{llncs}

\usepackage[english,brazil]{babel}
\usepackage[utf8]{inputenc}
\usepackage[T1]{fontenc}
\usepackage{graphicx}
\usepackage{booktabs}
\usepackage{hyperref}
\usepackage{enumitem}
\usepackage{listings}
\usepackage{xcolor}

\graphicspath{{./figures/}}

\lstset{
  basicstyle=\ttfamily\small,
  breaklines=true,
  frame=single,
  columns=fullflexible
}

\title{Temas de Pesquisa em Mobilidade Urbana\\Consolidação de Quatro Trabalhos Incompletos}
\titlerunning{Temas de Pesquisa em Mobilidade Urbana}
\author{Material consolidado para organização interna}
\authorrunning{Organização interna}
\institute{CEFET/RJ}
\date{}

\begin{document}
\maketitle

\begin{abstract}
Este documento consolida quatro trabalhos incompletos sobre mobilidade urbana em uma estrutura única de temas de pesquisa. O material foi reorganizado para separar, em cada tema, a introdução, o referencial teórico, a proposta de metodologia, os resultados parciais e os códigos associados. Além disso, foi criado um capítulo inicial comum com o procedimento de ETL utilizado para obtenção e preparação do dataset base empregado pelos estudos.
\end{abstract}

\section{Capítulo 0: ETL Comum}

\subsection{Introdução}
Os quatro trabalhos compartilham a mesma base de dados de mobilidade de ônibus do Rio de Janeiro. Por isso, o processo de extração e preparação inicial foi tratado como uma etapa comum. O objetivo desse capítulo é registrar como o dataset pode ser buscado, lido e salvo para uso posterior nos experimentos.

\subsection{Referencial}
O ponto de partida do conjunto de estudos é o uso de dados abertos de mobilidade urbana, em especial registros de GPS de ônibus combinados com atributos operacionais e contextuais. A documentação deixada na pasta original aponta para a fonte pública do dataset e para um fluxo simples de leitura de arquivos Parquet seguido da persistência em \texttt{RData}.

\subsection{Proposta de Metodologia}
O fluxo comum foi organizado em duas etapas. Primeiro, um script em Python lê arquivos Parquet remotos. Em seguida, um script em R orquestra o carregamento dos cinco fragmentos do dia analisado e salva um arquivo único para os estudos subsequentes.

\begin{lstlisting}[language=Python,caption={Leitura simples de arquivos Parquet em Python}]
import pandas as pd

def read_parquet(filename):
    df = pd.read_parquet(filename)
    return df
\end{lstlisting}

\begin{lstlisting}[language=R,caption={Orquestração da carga do dataset em R}]
library(reticulate)

source_python("read-parquet.py")

data <- read_parquet("https://regulus.eic.cefet-rj.br/data/mobility/G1-2023-02-21-A.parquet")
data_tmp <- read_parquet("https://regulus.eic.cefet-rj.br/data/mobility/G1-2023-02-21-B.parquet")
data <- merge(data, data_tmp)
\end{lstlisting}

\subsection{Resultados Parciais}
Como resultado parcial, essa etapa produz uma base consolidada pronta para filtragens por linha, veículo, horário, velocidade e localização. Esse ETL serve como dependência explícita para os quatro temas.

\subsection{Códigos Associados}
\begin{itemize}[leftmargin=*]
\item \texttt{codigos/00-etl/read-parquet.py}
\item \texttt{codigos/00-etl/read-parquet.R}
\item \texttt{codigos/00-etl/documentacao.txt}
\end{itemize}

\section{Tema 1: Identificação de Congestionamentos com DBSCAN}

\subsection{Introdução}
Este tema investiga congestionamentos em linhas de ônibus a partir de dados de GPS, velocidade e chuva. A proposta parte do entendimento de que atrasos e retenções não aparecem apenas como eventos explícitos, mas também como padrões espaço-temporais recorrentes distribuídos pela cidade.

\subsection{Referencial Teórico}
O núcleo teórico combina três frentes. A primeira é mobilidade urbana, com foco em fluidez do transporte coletivo. A segunda é aprendizado não supervisionado, especialmente agrupamento por densidade. A terceira é o algoritmo DBSCAN, escolhido por identificar agrupamentos sem exigir a quantidade de clusters previamente definida e por lidar razoavelmente bem com ruído.

\subsection{Proposta de Metodologia}
A metodologia definida no material original segue quatro passos:
\begin{enumerate}[leftmargin=*]
\item Conversão do timestamp para hora decimal e filtragem do horário de pico.
\item Seleção de ônibus com velocidade muito baixa e exclusão de casos de estacionamento.
\item Normalização de latitude, longitude, horário e volume de chuva.
\item Aplicação de DBSCAN para identificar agrupamentos de pontos com comportamento compatível com congestionamento.
\end{enumerate}

O script principal também produz recortes por região administrativa, permitindo observar padrões no Centro, em Campo Grande e em outras áreas.

\subsection{Resultados Parciais}
Os resultados parciais do trabalho indicam que o método consegue destacar regiões com concentração de eventos de baixa velocidade no horário de pico. O relatório original aponta maior incidência em áreas como Campo Grande e Centro, mas ainda reconhece limitações importantes, sobretudo o uso de um recorte temporal curto e a dificuldade em separar congestionamento recorrente de impacto climático pontual.

\subsection{Códigos Associados}
\begin{itemize}[leftmargin=*]
\item \texttt{codigos/tema-01-dbscan/analise-congestionamento.R}
\end{itemize}

\subsection{Possíveis Expansões para Artigo Científico}
\begin{itemize}[leftmargin=*]
\item Ampliar a janela temporal para múltiplos dias, semanas e períodos chuvosos e não chuvosos.
\item Comparar o DBSCAN com métodos alternativos, como HDBSCAN, OPTICS e abordagens supervisionadas.
\item Construir uma métrica formal de severidade de congestionamento combinando velocidade, duração, localização e chuva.
\item Validar os agrupamentos com fontes externas, como dados viários, notícias de trânsito ou registros oficiais.
\item Explorar predição de congestionamentos a partir de séries históricas e variáveis contextuais.
\end{itemize}

\section{Tema 2: Classificação do Sentido de Trajeto com Produto Escalar}

\subsection{Introdução}
Este tema aborda a ausência do campo de direção de viagem em bases de mobilidade. O problema central é distinguir, para cada observação de GPS, se o ônibus está em trajeto de ida ou de volta, mesmo quando os dois sentidos compartilham trechos fisicamente sobrepostos.

\subsection{Referencial Teórico}
O referencial do tema é formado por dados de mobilidade urbana, georreferenciamento, cálculo geodésico e álgebra vetorial. A fórmula de Haversine é utilizada para medir proximidade entre observações e terminais, enquanto o produto escalar é empregado para comparar o deslocamento instantâneo do veículo com um vetor de referência definido pelos terminais da linha.

\subsection{Proposta de Metodologia}
A proposta metodológica foi consolidada em um pipeline computacional:
\begin{enumerate}[leftmargin=*]
\item Definição manual de dois terminais por linha.
\item Cálculo das distâncias de cada ponto até os terminais.
\item Rotulagem da localização como próximo ao terminal A, ao terminal B ou em trânsito.
\item Cálculo do produto escalar entre o vetor de deslocamento e o vetor de referência A$\rightarrow$B.
\item Segmentação do trajeto com base na alternância entre terminais.
\item Classificação final do segmento por soma acumulada dos sinais direcionais.
\end{enumerate}

Essa solução procura reduzir oscilações locais causadas por curvas, desvios ou pequenas inversões momentâneas de direção.

\subsection{Resultados Parciais}
Os resultados parciais mostram que a solução funciona bem para trajetos completos das linhas testadas e melhora em relação ao uso isolado do produto escalar. Ainda assim, permanecem fragilidades em cenários de garagem, inconsistência no campo de linha e escolha inadequada do buffer de proximidade.

\subsection{Códigos Associados}
\begin{itemize}[leftmargin=*]
\item \texttt{codigos/tema-02-direcao/classify.py}
\item \texttt{codigos/tema-02-direcao/cumsum.py}
\item \texttt{codigos/tema-02-direcao/cumsum-tendencia.py}
\item \texttt{codigos/tema-02-direcao/produto-escalar.py}
\item \texttt{codigos/tema-02-direcao/read-data.R}
\end{itemize}

\subsection{Possíveis Expansões para Artigo Científico}
\begin{itemize}[leftmargin=*]
\item Automatizar a descoberta dos terminais, reduzindo dependência de parametrização manual.
\item Comparar a abordagem vetorial com modelos supervisionados e métodos probabilísticos para inferência de direção.
\item Formalizar um protocolo de avaliação com rotulagem manual de amostras e métricas como acurácia, precisão e F1-score.
\item Tratar explicitamente cenários de garagem, desvios operacionais e inconsistências cadastrais da linha.
\item Escalar o método para um conjunto maior de linhas e testar sua robustez em trajetos com maior sobreposição espacial.
\end{itemize}

\section{Tema 3: Detecção de Anomalias em Séries Temporais de Velocidade}

\subsection{Introdução}
Este tema trata a velocidade dos ônibus como série temporal e busca detectar comportamentos anômalos. A motivação é identificar registros inconsistentes, períodos atípicos de operação ou padrões incomuns que possam sinalizar problemas de coleta ou de operação.

\subsection{Referencial Teórico}
O referencial combina análise de séries temporais, detecção de anomalias e preparação de dados em ambiente R. O relatório original menciona o uso do framework Harbinger para detecção de eventos temporais e da Daltoolbox como apoio ao fluxo analítico.

\subsection{Proposta de Metodologia}
O material consolidado mostra um fluxo com duas partes. A primeira organiza a base por linha, veículo e timestamp, transformando as observações em uma tabela temporal por ônibus. A segunda remove valores problemáticos e prepara a série para análise de anomalias por modelo estatístico. A linha 321 e o veículo B28519 aparecem como estudo de caso principal no relatório.

\subsection{Resultados Parciais}
Os resultados parciais sugerem que havia valores constantes e trechos suspeitos que comprometiam a leitura inicial dos dados. A revisão do dataset e a remoção desses comportamentos tornaram a análise mais coerente, permitindo um cenário mais favorável à aplicação do detector de anomalias. O estudo, porém, ainda se encontra em estágio exploratório e depende de refino metodológico.

\subsection{Códigos Associados}
\begin{itemize}[leftmargin=*]
\item \texttt{codigos/tema-03-anomalias/series.R}
\item \texttt{codigos/tema-03-anomalias/mobilidade-1.R}
\end{itemize}

\subsection{Possíveis Expansões para Artigo Científico}
\begin{itemize}[leftmargin=*]
\item Definir taxonomias de anomalias operacionais, como parada prolongada, velocidade incompatível e falha de sensor.
\item Comparar modelos ARIMA, detecção baseada em decomposição sazonal e métodos modernos de aprendizado de máquina.
\item Avaliar o impacto do pré-processamento sobre a taxa de detecção e sobre falsos positivos.
\item Expandir o estudo para múltiplas linhas e diferentes dias, investigando recorrência das anomalias.
\item Relacionar anomalias detectadas com eventos externos, como chuva, horários de pico e mudanças de operação.
\end{itemize}

\section{Tema 4: Classificação de Trajetórias com WiSARD}

\subsection{Introdução}
Este tema investiga a classificação de trajetórias de ônibus por linha a partir de redes neurais sem peso, com foco na arquitetura WiSARD. O objetivo é transformar séries de coordenadas geográficas em representações binárias adequadas ao classificador e comparar diferentes estratégias de codificação.

\subsection{Referencial Teórico}
O referencial teórico reúne mobilidade urbana, classificação de trajetórias, binarização de dados geoespaciais, KernelCanvas, grades uniformes, grades baseadas em KD-Tree, codificação por termômetro distributivo e a própria rede WiSARD. O material também discute hiperparâmetros como \texttt{addressSize}, \texttt{numberOfKernels}, \texttt{bitsByKernel} e \texttt{activationDegree}.

\subsection{Proposta de Metodologia}
O trabalho foi reorganizado em três blocos metodológicos:
\begin{enumerate}[leftmargin=*]
\item Preparação de trajetórias e transformação dos registros de GPS em imagens ou vetores binários.
\item Aplicação de diferentes técnicas de codificação, incluindo KernelCanvas e termômetro distributivo.
\item Treinamento e avaliação do classificador WiSARD para reconhecer linhas de ônibus e medir confiança da classificação.
\end{enumerate}

Essa formulação permite comparar precisão e custo computacional entre diferentes estratégias de representação.

\subsection{Resultados Parciais}
Os resultados parciais indicam melhor desempenho para a combinação grade uniforme + KernelCanvas, com métricas superiores às demais abordagens descritas no relatório original. Em contrapartida, o termômetro distributivo apresentou menor custo computacional na etapa de binarização, o que sugere utilidade em cenários com maior restrição temporal.

\subsection{Códigos Associados}
\begin{itemize}[leftmargin=*]
\item \texttt{codigos/tema-04-wisard/Codec.py}
\item \texttt{codigos/tema-04-wisard/wisard/}
\item \texttt{codigos/tema-04-wisard/KernelCanvas/}
\end{itemize}

\subsection{Possíveis Expansões para Artigo Científico}
\begin{itemize}[leftmargin=*]
\item Realizar comparação sistemática com classificadores tradicionais e profundos, como Random Forest, SVM, CNN e LSTM.
\item Investigar novas formas de codificação geoespacial e temporal além de KernelCanvas e termômetro distributivo.
\item Estudar sensibilidade dos hiperparâmetros da WiSARD com protocolo experimental mais amplo e reprodutível.
\item Avaliar generalização em diferentes subconjuntos de linhas, regiões e períodos temporais.
\item Explorar a confiança da classificação como sinal auxiliar para detecção de desvios de rota e comportamento atípico.
\end{itemize}

\section{Conclusão}
Os quatro trabalhos foram convertidos em temas de pesquisa com uma organização homogênea. A nova estrutura permite enxergar com clareza o que já existe em cada frente: motivação, base teórica, caminho metodológico, estágio atual dos resultados e scripts de apoio. Com isso, o conjunto fica mais adequado para continuidade acadêmica, reaproveitamento de código e refinamento futuro das propostas.

\end{document}
